















[...]

[...]

[...]


\[
    \epsilon(z) = \sum_{n = -\infty}^{\infty} \epsilon_n z^{n + 1},
\]
unless you have a constant function, if it is analytic there must be somewhere in the domain of analiticity in which the function is greater or higher than some value, and outside the domain of analiticity the function can diverge. This laurent series is the most general form of an analytic function, and it is convergent in the annulus defined by the inner and outer radius of convergence.

We can consider \(\epsilon_n(z)\) our basis for the expansion
\[
    \epsilon_n(z) = - \epsilon z^{n + 1},
\]
then the generators of the transformations on scalar fields \(\varphi(x)\) as
\[
    \hat{\ell}_n = - z^{n + 1} \partial_z,
\]
such that
\begin{enumerate}
    \item \(\hat{\ell}_0 = -z \partial_z\) is the generator of \textbf{dilatations},
    \item \(\hat{\ell}_{-1} = -\partial_z\) is the generator of \textbf{translations},
    \item \(\hat{\ell}_1 = -z^2 \partial_z\) is the generator of \textbf{special conformal transformations}.
\end{enumerate}

The dilatation or translation is in \(z\), but we need the \(\overline{z}\) direction, such that we need to define
\[
    \overline{\ell}_n = - \overline{z}^{n + 1} \partial_{\overline{z}}.
\]
So we have
\begin{enumerate}
    \item \(\overline{\ell}_0 = -\overline{z} \partial_{\overline{z}}\) is the generator of \textbf{dilatations},
    \item \(\overline{\ell}_{-1} = -\partial_{\overline{z}}\) is the generator of \textbf{translations},
    \item \(\overline{\ell}_1 = -\overline{z}^2 \partial_{\overline{z}}\) is the generator of \textbf{special conformal transformations}.
\end{enumerate}

Thus if we were to write the four momentum operator in the complex coordinates, we would have
\[
    P = \begin{pmatrix}
        E \\
        p
    \end{pmatrix} = \begin{pmatrix}
        \ell_{-1} + \overline{\ell}_{-1} \\
        i(\ell_{-1} - \overline{\ell}_{-1})
    \end{pmatrix},
\]
where \(E\) is the energy and \(p\) is the momentum. The dilatation operator would be
\[
    D = \hat{\ell}_0 + \overline{\ell}_0,
\]

[...]

If you reason in higher dimensions, the special conformal group is six-dimensional, since
\[
    D=2 \implies \frac{(D+1)(D+2)}{2} = 6.
\]

Imagine to apply the commutator \(\left[ \ell_n, \, \ell_m \right]\) on a test function \(f(z)\), we can find
\[
    \left[ \ell_n, \, \ell_m \right] = \left( (n - m) z^{n + m + 1} \partial_z \right) = (n - m) \ell_{n + m},
\]
which is the Witt algebra, the algebra of the generators of conformal transformations in two dimensions. The same holds for the barred generators \(\overline{\ell}_n\). The Witt algebra is an infinite-dimensional Lie algebra, since \(n\) and \(m\) can take any integer value and be applied on any test function \(f(z)\) in the domain of analyticity. The Witt algebra is the classical version of the Virasoro algebra, which is the quantum version of the conformal algebra in two dimensions. Visaroro was a young phd in Turin when he discovered the algebra, after is supervisor (which already had found the Mobius transformations we are going to introduce later) told him to find the algebra of the generators of conformal transformations in two dimensions, and he found this infinite-dimensional algebra.

This generators, with \(n = -1,0,1\) (barred and not), are six generators which form the algebra of the \(2\times 2\) special linear group \(\mathfrak{sl}(2, \, \mathbb{C})\), which is the double cover of the Lorentz group \(SO(3,1)\). This is the conformal group in two dimensions. It is isomorphic to \(SU(2)\oplus SU(2)\).

We have an infinite number of conformal transformations, since they are all the analytic functions, so how can it be six dimensional? We can map the coordintes of the plane to \(z=e^{w}\), which inverse gives the logarithm of \(z\): we are mapping the plane to a strip (given by the riemann surface of the logarithm), and with periodic boundary conditions we get a cylinder. The conformal transformations on the plane are mapped to conformal transformations on the cylinder, and the six generators of the conformal group are mapped to the six generators of the conformal group on the cylinder. very important (it is a way of compactifying the plane).

Now we can ask which are the transformations starting from a plane and ending in a plane, i.e. which are the transformations that do not change the point at infinity. These are the Mobius transformations, which are given by
\[
    w = \frac{az + b}{cz + d}, \quad ad - bc = 1,
\]
which has six parameters. Thus global conformal transformations are given by the Mobius transformations, which are a subgroup of the infinite-dimensional group of local conformal transformations.

    [...]

We can now infer someting about the fields: after the remapping
\[
    \begin{dcases}
        z \\
        \overline{z}
    \end{dcases} \mapsto \begin{dcases}
        w = f(z) \\
        \overline{w} = \overline{f}(\overline{z})
    \end{dcases},
\]
we can find fields that transforms as
\[
    \dots
\]
which are called \textbf{primary fields}, and they are the building blocks of the theory. We can also understand quasi-primary fields, which are fields that transforms as primary fields under the global conformal transformations (Mobius transformations), but not under the local conformal transformations. There are also other fields which do not have tensorial stucture under the conformal transformations, but they are not important for us. Some of the secondary fields happens to be quasi-primary, but not all of them. The primary fields are the most important fields in the theory and we will distiguish them from the secondary fields, which are obtained by applying the generators of the conformal transformations on the primary fields.

For example, the energy-momentum tensor \(T(z)\) is secondary field, and it would seem lika a bad news since it is the most important field in the theory, but we will see the consequences.

If we compute tge derivatives of \(w\) and \(\overline{w}\) with respect to \(z\) and \(\overline{z}\), we can find
\[
    \dots
\]
and now expanding
\[
    \phi(z, \, \overline{z}) = \phi(w - \epsilon, \, \overline{w} - \overline{\epsilon}) = \phi(w, \, \overline{w}) - \epsilon \partial_w \phi(w, \, \overline{w}) - \overline{\epsilon} \partial_{\overline{w}} \phi(w, \, \overline{w}) + \dots
\]
such that we can find the variation of the field \(\phi\) under the conformal transformation
\[
    \delta \phi = - \epsilon \partial_z \phi - \overline{\epsilon} \partial_{\overline{z}} \phi - h \partial_z \epsilon \phi - \overline{h} \partial_{\overline{z}} \overline{\epsilon} \phi,
\]
where we can see how the two variables are completely decoupled, and the transformation of the field is given by the sum of the transformation in \(z\) and the transformation in \(\overline{z}\).

dilatations\dots
\[
    D = \ell_0 + \overline{\ell}_0 = h + \overline{h} = \Delta,
\]

rotations\dots
\[
    L = \ell_0 - \overline{\ell}_0 = h - \overline{h} = s,
\]

such that we can characterize the fields by their scaling dimension \(\Delta\) and their spin \(s\) (or equivalently by \(h\) and \(\overline{h}\)). A scalar field is spinless, so \(s=0\) and \(h = \overline{h}\). A chiral field is a field that depends only on \(z\) or only on \(\overline{z}\), so it has either \(h=0\) or \(\overline{h}=0\).